\documentclass[10pt,a4paper]{amsart}
\usepackage[margin=19mm]{geometry}
\usepackage{amsmath,amssymb,mathtools}
\usepackage[scr=rsfs,cal=euler]{mathalfa}
\usepackage{xfrac}
\usepackage[unicode,hidelinks,bookmarks=false]{hyperref}
\newcommand{\ie}{i.e.\ }
\newcommand{\cf}{cf.\ }
\newcommand{\resp}{respectively\ }
\newcommand{\Subsection}{\subsection}
\providecommand{\nobreakdash}{\nobreak\hbox{-}}
\setlength{\emergencystretch}{2em}
\multlinegap0pt
\usepackage{xcolor}
\usepackage{amssymb}
\newtheorem{theorem}{Theorem}
\providecommand{\paragraph}{}
\title{Optimal Guarantees for Auditing R\'enyi Differentially Private Machine Learning}
\author{Benjamin D. Kim, Lav R. Varshney, Daniel Alabi}
\date{Source: arXiv:2605.21938v1 (CC BY 4.0)}
\begin{document}
\maketitle

\noindent\emph{Source and licence.} Optimal Guarantees for Auditing R\'enyi Differentially Private Machine Learning. Version arXiv:2605.21938v1 (CC BY 4.0), distributed by arXiv under the Creative Commons Attribution 4.0 International licence, http://creativecommons.org/licenses/by/4.0/, as recorded in the arXiv licence metadata for this exact version and shown on the abstract page https://arxiv.org/abs/2605.21938v1. Full licence notice: \texttt{licenses/arxiv-2605.21938-CC-BY-4.0.txt}.\par\smallskip

\noindent\textit{Editorial scope.} Counted unit: the article's theorem thm:finite-sample-dv-renyi-app (source0.tex lines 1413--1487). The article's complete original proof of it is delivered with it (source0.tex lines 1489--1602, one \texttt{proof} environment of 114 lines). No mathematics is written, repaired or re-derived: the statement and the proof are the article's own lines, in the article's own notation, and no retained line is substituted or re-flowed.  No further source-local context is needed: every cross-reference of every retained line resolves inside the two delivered spans.  Nothing was omitted from any retained span, so the package carries no omission record.  No retained line contains a citation, so the wrapper carries no bibliography and no bibliography entry.  No retained line contains a figure environment, an image-inclusion command or drawing code, so no image is delivered and the package declares no asset dependency.  Editorial formatting only, in addition: an AMS \texttt{amsart} preamble that restates the article's own declarations and macros used by the retained lines, namely the environments \texttt{theorem} on separate counters, together with the standard packages those lines need.  The retained environments print in this document as Theorem 1 in order of appearance, whereas the article prints the same items with the numbering of its own full document.  The complete native source of the article as distributed in the arXiv e-print for this exact version is bundled unchanged as \texttt{sources/201100/source0.tex}.
\begin{theorem}[Consistency and finite-sample confidence interval for DV R\'enyi estimators]
\label{thm:finite-sample-dv-renyi-app}
Let $\alpha\in\mathbb{R}_{> 0}\setminus\{1\}$.  
Let $P,Q$ be probability distributions on a measurable space $\Omega$.
Let $\{X_i\}_{i=1}^n \sim Q$ and $\{Y_i\}_{i=1}^n \sim P$ be independent samples.

Let $\Theta \subset \mathbb{R}^d$ be a parameter set with $\|\theta\|\le K$.
Assume the critic family $\{T_\theta:\Omega\to\mathbb{R}\}_{\theta\in\Theta}$ satisfies:
\begin{enumerate}
\item (\textbf{Uniform boundedness}) $\sup_{\theta,z}|T_\theta(z)| \le M$.
\item (\textbf{Lipschitz parameterization}) For all $z\in\Omega$,
\[
|T_\theta(z)-T_{\theta'}(z)| \le L \|\theta-\theta'\| .
\]
\end{enumerate}

Define the population DV R\'enyi functional
\[
V(\theta)
:=\frac{1}{\alpha-1}\log \mathbb{E}_{Q}\!\left[e^{(\alpha-1)T_\theta(X)}\right]
-\frac{1}{\alpha}\log \mathbb{E}_{P}\!\left[e^{\alpha T_\theta(Y)}\right],
\]
and its empirical estimator
\[
\widehat V_n(\theta)
:=\frac{1}{\alpha-1}\log\!\left(\frac{1}{n}\sum_{i=1}^n e^{(\alpha-1)T_\theta(X_i)}\right)
-\frac{1}{\alpha}\log\!\left(\frac{1}{n}\sum_{i=1}^n e^{\alpha T_\theta(Y_i)}\right).
\]

Define
\[
R_\alpha^{\Theta}(Q\|P):=\sup_{\theta\in\Theta} V(\theta),
\qquad
\widehat R_{\alpha,n}^{\Theta}(Q\|P):=\sup_{\theta\in\Theta} \widehat V_n(\theta).
\]

Then for any $\delta\in(0,1)$, with probability at least $1-\delta$,
\[
\big|\widehat R_{\alpha,n}^{\Theta}(Q\|P)-R_\alpha^{\Theta}(Q\|P)\big|
\;\le\;
\varepsilon_n(\delta),
\]
where
\[
\varepsilon_n(\delta)
=
C_{\alpha,M}
\left(
\sqrt{\frac{d\log\!\left(\frac{K}{\eta}\right)+\log(1/\delta)}{n}}
+
\eta
\right),
\]
for any $\eta>0$, and
\[
C_{\alpha,M}
=
O\!\left(
e^{2|\alpha|M}\max\Big\{\tfrac{1}{|\alpha-1|},\tfrac{1}{|\alpha|}\Big\}
\right).
\]

In particular, choosing $\eta = O(n^{-1/2})$ yields a valid $(1-\delta)$
finite-sample confidence interval
\[
\Big[
\widehat R_{\alpha,n}^{\Theta}(Q\|P)-\varepsilon_n(\delta),
\;
\widehat R_{\alpha,n}^{\Theta}(Q\|P)+\varepsilon_n(\delta)
\Big].
\]
In addition,
the estimator $\widehat R_{\alpha,n}^{\Theta}(Q\|P)$ is consistent in the following sense:
$\big|\widehat R_{\alpha,n}^{\Theta}(Q\|P)-R_\alpha^{\Theta}(Q\|P)\big|\rightarrow 0$ as $n\rightarrow\infty$.
\end{theorem}

\begin{proof}
The proof for the finite-sample bound proceeds in four steps.

\paragraph{Step 1: Concentration for fixed $\theta$.}
Define
\[
A_\theta := \mathbb{E}_{Q}[e^{(\alpha-1)T_\theta(X)}],
\qquad
\widehat A_\theta := \frac{1}{n}\sum_{i=1}^n e^{(\alpha-1)T_\theta(X_i)}.
\]
Since $|T_\theta|\le M$,
\[
e^{-|\alpha-1|M} \le e^{(\alpha-1)T_\theta(X)} \le e^{|\alpha-1|M}.
\]
Thus $\widehat A_\theta$ is an average of i.i.d.\ bounded random variables.
Hoeffding's inequality implies that for all $t>0$,
\[
\mathbb{P}\left(|\widehat A_\theta-A_\theta|\ge t\right)
\le
2\exp\left(-\frac{n t^2}{2e^{2|\alpha-1|M}}\right).
\]

Similarly, define
\[
B_\theta := \mathbb{E}_{P}[e^{\alpha T_\theta(Y)}],
\qquad
\widehat B_\theta := \frac{1}{n}\sum_{i=1}^n e^{\alpha T_\theta(Y_i)},
\]
and obtain
\[
\mathbb{P}\left(|\widehat B_\theta-B_\theta|\ge t\right)
\le
2\exp\left(-\frac{n t^2}{2e^{2|\alpha|M}}\right).
\]

\paragraph{Step 2: From moment errors to log errors.}
Since $A_\theta,\widehat A_\theta \ge e^{-|\alpha-1|M}$,
the logarithm is Lipschitz on this interval with constant $e^{|\alpha-1|M}$.
Thus
\[
|\log \widehat A_\theta-\log A_\theta|
\le
e^{|\alpha-1|M}|\widehat A_\theta-A_\theta|.
\]
An analogous bound holds for $\log \widehat B_\theta$.

Combining both terms,
\[
|\widehat V_n(\theta)-V(\theta)|
\le
\frac{e^{|\alpha-1|M}}{|\alpha-1|}|\widehat A_\theta-A_\theta|
+
\frac{e^{|\alpha|M}}{|\alpha|}|\widehat B_\theta-B_\theta|.
\]
Therefore there exists a constant $c>0$ such that
\[
\mathbb{P}\left(|\widehat V_n(\theta)-V(\theta)|\ge u\right)
\le
4\exp\left(-c\,n\,u^2 e^{-4|\alpha|M}\right).
\]

\paragraph{Step 3: Uniformization over $\Theta$ via covering numbers.}
Let $\mathcal{N}_\eta$ be an $\eta$-net of $\Theta$ in Euclidean norm.
Since $\Theta\subseteq B(0,K)\subset\mathbb{R}^d$,
\[
|\mathcal{N}_\eta| \le \left(\frac{3K}{\eta}\right)^d.
\]
Applying a union bound,
\[
\mathbb{P}\left(
\max_{\vartheta\in\mathcal{N}_\eta}
|\widehat V_n(\vartheta)-V(\vartheta)|
\ge u
\right)
\le
|\mathcal{N}_\eta|\cdot
4\exp\left(-c\,n\,u^2 e^{-4|\alpha|M}\right).
\]
Choosing
\[
u
=
C_{\alpha,M}
\sqrt{\frac{d\log(K/\eta)+\log(1/\delta)}{n}}
\]
ensures the above probability is at most $\delta/2$.

\paragraph{Step 4: Extension from the net to all parameters.}
The Lipschitz assumption implies that both $V(\theta)$ and $\widehat V_n(\theta)$
are Lipschitz in $\theta$ with constant
$O(|\alpha|Le^{2|\alpha|M})$.
Hence for any $\theta\in\Theta$ and nearest net point $\vartheta\in\mathcal{N}_\eta$,
\[
|\widehat V_n(\theta)-V(\theta)|
\le
|\widehat V_n(\vartheta)-V(\vartheta)| + O(\eta).
\]
Taking the supremum over $\theta$ yields
\[
\sup_{\theta\in\Theta}|\widehat V_n(\theta)-V(\theta)|
\le
\max_{\vartheta\in\mathcal{N}_\eta}|\widehat V_n(\vartheta)-V(\vartheta)| + O(\eta).
\]

\paragraph{Step 5: Confidence interval for the supremum.}
Finally, for any functions $f,g$,
\[
|\sup f - \sup g| \le \sup |f-g|.
\]
Applying this inequality completes the proof of the finite-sample bound.

The consistency statement follows from the finite-sample bound.

\end{proof}

\end{document}
